problem:  
Let \((M^n, g)\) be a complete, non-collapsed Riemannian manifold with nonnegative Ricci curvature, and let \(\tilde M\) denote its universal cover. For each sequence \(R_i \to \infty\), consider the equivariant rescalings \((\tilde M, R_i^{-2} g, \tilde p, \Gamma)\), where \(\Gamma = \pi_1(M)\) acts by deck transformations, and let \((X_\infty, d_\infty, p_\infty, G_\infty)\) be a limit in the *equivariant Gromov–Hausdorff sense*.  

By Cheeger–Colding–Naber theory, \(X_\infty\) is a metric cone \(C(Z)\) that splits as \(\mathbb{R}^{k_\infty} \times C(Z_0)\), and \(G_\infty\) acts by isometries on \(X_\infty\).  
Let \(t(G_\infty) \subset G_\infty\) be the **maximal linear translation subgroup**—the connected, abelian subgroup of isometries acting as pure translations on the full Euclidean factor \(\mathbb{R}^{k_\infty}\).  

**Open problem:**  
Does the rank of the translation subgroup in the equivariant asymptotic cone coincide with the maximal Euclidean splitting dimension, i.e.
\[
\dim t(G_\infty) = k_\infty?
\]
Equivalently, do all Euclidean directions appearing in the asymptotic cone of \(\tilde M\) arise from the limiting translation action of deck transformations under the equivariant Gromov–Hausdorff convergence?  
If not always, under what geometric hypotheses (e.g., maximal volume growth or uniqueness of tangent cones at infinity) does the equality hold?

Why is it a "good" problem:  
This formulation now rests entirely on rigorously defined objects within the existing framework of equivariant Gromov–Hausdorff limits. The translation subgroup \(t(G_\infty)\) is the canonical asymptotic remnant of the deck transformation action, making the problem mathematically meaningful yet open.  

The question probes an *equivariant correspondence* between algebraic and geometric asymptotics: whether the linear part of the limit group action accounts fully for all Euclidean factors in metric cone splitting. This would unify the analytic structure of Ricci-limit spaces with the algebraic structure of the isometry group’s limit, extending Cheeger–Gromoll’s local splitting to an asymptotic, group-theoretic equality. It precisely targets an unproven dictionary between **large-scale geometry** (Euclidean directions in tangent cones) and **group asymptotics** (translation ranks of fundamental group actions).  

If established, such an equality would create a fundamental new bridge connecting Ricci-limit geometry, group actions in asymptotic cones, and the algebraic decomposition theory of virtually nilpotent fundamental groups. The problem’s phrasing is concrete and minimal—comparing two ranks—but its implications reach deeply into the geometry of nonnegative Ricci spaces and the structure of their fundamental groups.